\section{Практичне дослідження цифрового підпису}
\subsection{Вибір криптографічного алгоритму та бібліотеки}

Для практичного дослідження обрано алгоритм Ed25519, який належить до сімейства алгоритмів 
цифрового підпису на основі еліптичних кривих. Він забезпечує створення та перевірку цифрових 
підписів і характеризується високою швидкістю роботи та відносно невеликим розміром ключів.

Для реалізації використано бібліотеку \texttt{cryptography} мови Python, яка підтримує роботу з Ed25519. 
Вона надає готові засоби для генерації ключів, створення цифрового підпису та його перевірки, 
тому самостійна реалізація математичних операцій алгоритму не потрібна \cite{RFC8032,cryptography}.


\subsection{Генерація пари ключів}

На першому етапі генерується пара ключів: закритий ключ та відповідний йому відкритий ключ. 
Закритий ключ використовується для створення цифрового підпису, а відкритий — для його перевірки.

\begin{lstlisting}[language=Python]
from cryptography.hazmat.primitives.asymmetric.ed25519 import Ed25519PrivateKey
from cryptography.exceptions import InvalidSignature

private_key = Ed25519PrivateKey.generate()
public_key = private_key.public_key()

print("1. Key pair generated.")
\end{lstlisting}

Результат виконання:

\begin{lstlisting}
1. Key pair generated.
\end{lstlisting}


\subsection{Створення тестового повідомлення}

Для проведення експерименту створюється тестове повідомлення. Оскільки криптографічні операції виконуються 
над байтовими даними, текст перетворюється за допомогою методу \texttt{encode()}.

\begin{lstlisting}[language=Python]
message = "This is a test message for digital signature verification."
message = message.encode("utf-8")

print("\n2. Test message:")
print(message.decode("utf-8"))
\end{lstlisting}

Результат виконання:

\begin{lstlisting}
2. Test message:
This is a test message for digital signature verification.
\end{lstlisting}


\subsection{Формування цифрового підпису}

Для створення цифрового підпису використовується закритий ключ та підготовлене повідомлення. Результатом є 
послідовність байтів, яка надалі використовуватиметься для перевірки автентичності повідомлення.

\begin{lstlisting}[language=Python]
signature = private_key.sign(message)

print("\n3. Digital signature created.")
print("Signature:", signature.hex())
\end{lstlisting}

Результат виконання:

\begin{lstlisting}
3. Digital signature created.
Signature: 3c3a30bd02c73ba3daf57323007d6dcd828f2d51b1f0a01df2a63f90ab7d95d6236f0ea1c
9b9d5ab7aa5320ada4a30b3b68f57deb5c78f317b7e9fa546937904
\end{lstlisting}

Значення підпису залежить від згенерованої пари ключів, тому під час кожного запуску програми воно буде іншим.


\subsection{Перевірка цифрового підпису}

На цьому етапі виконується перевірка створеного підпису за допомогою відкритого ключа. Якщо підпис відповідає 
повідомленню та відкритому ключу, перевірка завершується успішно. Інакше бібліотека генерує виняток \texttt{InvalidSignature}.

\begin{lstlisting}[language=Python]
print("\n4. Verification of the original message:")

try:
    public_key.verify(signature, message)
    print("Signature is valid.")
except InvalidSignature:
    print("Signature is invalid.")
\end{lstlisting}

Результат виконання:

\begin{lstlisting}
4. Verification of the original message:
Signature is valid.
\end{lstlisting}


\subsection{Перевірка підпису після зміни повідомлення}

Для перевірки цілісності повідомлення змінюється його вміст, після чого виконується повторна перевірка початкового 
цифрового підпису. Очікується, що перевірка завершиться невдало, оскільки підпис був створений для іншого повідомлення.

\begin{lstlisting}[language=Python]
modified_message = "This is a modified test message.".encode("utf-8")

print("\n5. Verification of the modified message:")

try:
    public_key.verify(signature, modified_message)
    print("Signature is valid.")
except InvalidSignature:
    print("Signature is invalid: message was modified.")
\end{lstlisting}

Результат виконання:

\begin{lstlisting}
5. Verification of the modified message:
Signature is invalid: message was modified.
\end{lstlisting}


\subsection{Перевірка підпису з використанням іншого відкритого ключа}

На останньому етапі генерується інша пара ключів. Для перевірки початкового підпису використовується новий відкритий ключ, 
який не відповідає закритому ключу, за допомогою якого було створено підпис.

\begin{lstlisting}[language=Python]
other_private_key = Ed25519PrivateKey.generate()
other_public_key = other_private_key.public_key()

print("\n6. Verification with another public key:")

try:
    other_public_key.verify(signature, message)
    print("Signature is valid.")
except InvalidSignature:
    print("Signature is invalid: another public key was used.")
\end{lstlisting}

Результат виконання:

\begin{lstlisting}
6. Verification with another public key:
Signature is invalid: another public key was used.
\end{lstlisting}


\subsection{Аналіз результатів}

У результаті експерименту було реалізовано повний цикл роботи з цифровим підписом за допомогою алгоритму Ed25519.

Під час перевірки початкового повідомлення з використанням відповідного відкритого ключа підпис було успішно підтверджено. 
Після зміни повідомлення перевірка завершилася невдало, оскільки цифровий підпис не відповідає зміненим даним. 
Аналогічно, використання іншого відкритого ключа також призвело до невдалої перевірки.

Отримані результати демонструють основні властивості цифрового підпису:

\begin{enumerate}
    \item Цілісність — дозволяє виявити зміну повідомлення після створення підпису.
    \item Автентичність — дає змогу перевірити, чи відповідає підпис закритому ключу, по\-в'язаному з використаним відкритим ключем.
    \item Невідмовність — за умови надійного зберігання закритого ключа та підтвердження його належності конкретному власнику, 
    підпис може використовуватися як доказ авторства повідомлення.
\end{enumerate}

Таким чином, експеримент підтвердив, що цифровий підпис дозволяє контролювати цілісність повідомлень та перевіряти їхнє походження. 
При цьому він не забезпечує конфіденційності, оскільки саме повідомлення залишається доступним для читання.
